
\documentclass{screenplay}[2012/06/30]
\usepackage{verbatim}
\input{Template/commands}

% LEER INSTRUCTIVO.TXT PRIMERO
 
\begin{document}

\begin{comment}
Aquí se inserta la portada. Eso hace el comando \input. 
Si quisieras una portada con sólo el título: 
\input{Template\cover_clean} 

Si quisieras una portada normal: \input{Template/cover}
Si quisieras una portada normal pero sin info de contacto: \input{Template/cover_clean_autor}
\end{comment}

%Portada:
\input{Template/cover}

%Ignorar
\thispagestyle{empty}



%  AQUÍ EMPIEZA TU GUION
\fadein

\intslug[Día]{Casa}

JUAN (23) está parado en su sala, de mal humor... Lorem ipsum dolor sit amet, consectetur adipiscing elit. Suspdisse in mi porttitor, efficitur quam non.

\begin{dialogue}[]{Juan}
    No lo puedo creer
\end{dialogue}

Juan se acerca y bla bla bla

\begin{dialogue}[Enojada]{María}
    ¿Cómo que no?
\end{dialogue}

\begin{dialogue}[]{Juan \contd}
    Neta no
\end{dialogue}

\corte

\intslug[Día]{Oficina}
Otra escena con una descripción nueva

\final
\end{document}
